\preface[Introduction]
\label{intro:book}
\addcontentsline{toc}{chapter}{Introduction}

A singular potential is more than a function with an infinite value at a few points. Its derivatives can be measures, and those measures can charge sets invisible to Lebesgue integration. In several complex variables, the same potential also controls the integrability of holomorphic functions. Neither its logarithmic growth nor these integrability conditions determine all of its nonlinear mass. Understanding the relations between these objects requires a foundation in analysis before one can study singularities globally.

This book supplies that foundation for Part I of \emph{Singularities in global pluripotential theory} \cite{XiaSampleBook}. The starting point is classical potential theory; the destination is the analytic machinery behind non-pluripolar products, singularity envelopes, relative energy and the ample toric setting. The order follows the dependence of the proofs. In particular, weighted holomorphic estimates precede multiplier ideals, approximation precedes global positive classes, and the bounded Monge--Amp\`ere operator precedes every product involving unbounded potentials.

\section*{From derivatives to pointwise potentials}
Distributional differentiation makes sense for a locally integrable function, but forgets its values on a null set. Potential theory recovers those values from mean inequalities. A real distribution with positive Laplacian has a canonical subharmonic representative; harmonic replacement then controls maxima, boundary values and envelopes. The exceptional sets of these operations are measured by capacity rather than by Lebesgue measure. This is why measure theory, the Dirichlet problem, kernel energy and fine topology belong to the same foundation.

Plurisubharmonicity imposes the submean inequality on every complex line. A separate inequality on each line does not immediately give a uniform bound in a neighborhood: its constants may depend on the line. The Fourier and incidence estimates in \cref{chap:p06} close this gap. They allow the pointwise definition to be compared with positivity of the complex Hessian as a current. The resulting compactness and envelope theory concern the original pointwise functions, not only their almost-everywhere classes.

\section*{Two ways of measuring a singularity}
The first is local and holomorphic. Lelong numbers measure logarithmic concentration, while multiplier ideals record which holomorphic functions remain square integrable against an exponential weight. Proving that these ideals are coherent, and that integrability persists after a small increase of the weight, requires holomorphic extension. Weighted estimates therefore enter as constructions: they correct cutoff errors, prescribe values along slices, and produce the holomorphic functions used in approximation. Analytic singularities make the resulting growth explicit through finitely many holomorphic equations and, on a resolution, through divisors.

The second is nonlinear. For locally bounded plurisubharmonic functions, the Bedford--Taylor construction defines products of their curvature currents and proves continuity along decreasing sequences. For an unbounded function, one first truncates the potential and then discards the part of the product on its pole set. When the retained products have locally bounded mass, their limit is the non-pluripolar product. This operation can lose mass. Its relation to singularity order, and the conditions under which mass is preserved, are central to the global theory.

These two measurements need not agree. The construction in \cref{prop:diffuse-polar-cantor-measure,ex:diffuse-polar-zero-nonpluripolar-mass} gives an atom-free measure carried by a polar Cantor set and a potential on the projective line whose Lelong numbers vanish and whose multiplier ideals are all trivial, although its first-degree non-pluripolar mass is zero. This supplies the foundations of \cite[Example~6.1.3]{XiaSampleBook}. The example explains why neither logarithmic growth nor multiplier ideals can replace the mass theory developed later.

\section*{Global constraints and envelopes}
On a compact K\"ahler manifold, a smooth closed real $(1,1)$-form fixes a background cohomology class. Potentials vary the current within that class, so local positivity must be reconciled with a global constraint. Scalar elliptic theory supplies global potentials; compactness and approximation then describe pseudo-effective and big classes, including their boundary behavior. Singular Hermitian metrics express the same positivity in line-bundle language.

An envelope is the largest admissible potential below an obstacle. When an additional singularity constraint is imposed, it becomes a tool for selecting a representative of that constraint. Its Monge--Amp\`ere measure is controlled by contact with the obstacle, but this conclusion requires local replacement and precise limit arguments. Relative capacity and integration by parts then lead to energy, equations with prescribed singularities and Perron segments. The arbitrary-mass ceiling identity of \cite[Theorem~1.2]{PSWZ26} is proved in \cref{thm:arbitrary-mass-ceiling-envelope}; the zero-mass case is obtained from positive-mass perturbations with a uniform tail estimate. Convex duality and subgradient measures provide the parallel foundations needed on the open torus of a toric variety.

\section*{Prerequisites and reading order}
The reader is assumed to know undergraduate real analysis, including Lebesgue integration and its convergence theorems, linear algebra, point-set topology and elementary one-variable complex analysis. Distribution theory, the functional analysis used in weighted estimates, several-variable holomorphic tools, complex differential geometry and convex measure theory are developed here. The two parts form a continuous route: \cref{part:potential} establishes the linear and exceptional-set theory, and \cref{part:pluripotential} builds the complex local theory before passing to global singular potentials.

The appendix has a different role. It supplies the local analytic and geometric tools needed for multiplier ideals and modifications, and can be read when those tools first arise. The existence of analytic resolutions and marked-order-one principalizations is the explicit external interface in \cref{thm:analytic-resolution-input,thm:marked-order-one-input}. Their existence algorithms are not proved here; the uses of the resulting spaces in the analytic arguments are proved. Historical notes at the end identify sources and exact versions, with internal references connecting them to the mathematics.

The notation and normalizations follow \cite{XiaSampleBook}; they are declared in the conventions and where the relevant objects are introduced. This book develops the missing foundational arguments of its first part. Definitions and a limited number of shared constructions are repeated where a connected proof requires them, while the later classification and applications remain in the continuation.
